% This file was adapted from ICLR2022_conference.tex example provided for the ICLR conference
\documentclass{article} % For LaTeX2e
\usepackage{collas2025_conference,times}
\usepackage{easyReview}

\usepackage{mathtools}
\usepackage{subcaption}
\usepackage[dvipsnames]{xcolor}

\usepackage{algorithm}
\usepackage{algorithmic}

\usepackage{tikz}
\usepackage{amsmath, amssymb}
\usepackage{xcolor}

% Optional math commands from https://github.com/goodfeli/dlbook_notation.
\input{math_commands.tex}

% Please leave these options as they are
\usepackage{hyperref}
\hypersetup{
    colorlinks=true,
    linkcolor=red,
    filecolor=magenta,
    urlcolor=blue,
    citecolor=purple,
    pdftitle={Overleaf Example},
    pdfpagemode=FullScreen,
    }
\providecommand{\doi}{}
\renewcommand{\doi}[1]{\href{https://doi.org/#1}{\nolinkurl{doi:#1}}}




% \title{Group Empowerment}
\title{
Multi-Agent Empowerment \\ and Emergence of Complex Behavior in Groups
}

% Authors must not appear in the submitted version. They should be hidden
% as long as the \collasfinalcopy macro remains commented out below.
% Non-anonymous submissions will be rejected without review.

\author{Antiquus S.~Hippocampus, Natalia Cerebro  \thanks{ Use footnote for providing further information
about author (webpage, alternative address)---\emph{not} for acknowledging
funding agencies.  Funding acknowledgements go at the end of the paper.} \\
Department of Computer Science\\
Random University\\
Country \\
\texttt{\{hippo,brain\}@cs.random.edu} \\
\And % Use And to have authors side by side
Koala Learnus \& D. Q. ResNet  \\
Department of Computational Neuroscience \\
University of Random City \\
Another Country \\
\texttt{\{koala,net\}@random.rand} \\
\AND % Use AND to have authors block one under the other
Coauthor \\
Affiliation \\
Address \\
\texttt{email}
}

% The \author macro works with any number of authors. There are two commands
% used to separate the names and addresses of multiple authors: \And and \AND.
%
% Using \And between authors leaves it to \LaTeX{} to determine where to break
% the lines. Using \AND forces a linebreak at that point. So, if \LaTeX{}
% puts 3 of 4 authors names on the first line, and the last on the second
% line, try using \AND instead of \And before the third author name.

\newcommand{\fix}{\marginpar{FIX}}
\newcommand{\new}{\marginpar{NEW}}

% \collasfinalcopy % Uncomment for camera-ready version, but NOT for submission.

%\preprintcopy % Uncomment for the preprint version, but NOT for submission.

\begin{document}


\maketitle

\begin{abstract}
  Intrinsic motivations are receiving increasing attention, i.e. behavioral incentives that are not engineered, but emerge from the interaction of an agent with its surroundings. In this work we study the emergence of behaviors driven by one such incentive, empowerment, specifically in the context of more than one agent. We formulate a principled extension of empowerment to the multi-agent setting, and demonstrate its efficient calculation. We observe that this intrinsic motivation gives rise to characteristic modes of group-organization in two qualitatively distinct environments: a pair of agents coupled by a tendon, and a controllable Vicsek flock. This demonstrates the potential of intrinsic motivations such as empowerment to not just drive behavior for only individual agents but also higher levels of behavioral organization at scale.
\end{abstract}

\section{Introduction}

% \begin{figure}[H]
%     \centering
%     \includegraphics[width=\linewidth]{CoLLAs/Figures/Multi_User_Interference.pdf}
%     \caption{Multi-agent system as an interference channel with $N=3$ agents. Each agent ($n$) attempts to transmit its message (action sequence $U_{0:t-1}^{(n)}$) through the noisy channel to communicate with its receiver (the agent's future state $X_t^{(n)}$). The interference noise generated by the other channel users affects each agent's ability to communicate with its receiver.}
%     \label{fig:placeholder}
% \end{figure}


%% main question
Reinforcement learning remains the standard framework for training agents, but it relies on externally engineered rewards. These rewards are specific to particular tasks, hard to design canonically, and yield brittle task-specific solutions. Biological organisms, in contrast, rarely have well-delineated extrinsic tasks, yet they continuously learn and adapt. A growing body of work instead explores \emph{intrinsic motivations} \citep{tiomkin_2024} --- principles that generate behavior from an agent's interaction with its surroundings, without reference to external tasks. Intrinsically motivated behaviors are generic: they emerge from the structure of the agent--environment coupling rather than from a designer's choice of reward. Such generic behaviors can then be combined and refined into more complex task-specific behaviors; they can also equip an agent with a local incentive structure that cushions it against random or adversarial disturbances that would derail a naively task-optimized strategy \citep{tiomkin2025process}. This makes intrinsic motivation a natural substrate for asking what behaviors interacting agents produce on their own, without shared objectives --- the question we study here.

%% impose behaviors on agents via incentives

As autonomous agents become increasingly prevalent in virtual and physical environments, understanding what behaviors can emerge in agent groups from purely intrinsic motivation becomes an important question. We ask: when agents pursue only their own intrinsic motivation, with no explicit coordination or shared objective, what behaviors, if any, emerge from their interaction, and can we understand why? And how do those behaviors change when an agent instead directs that same capacity toward helping another?

%% what is empowerment, empowerment in multiagent cases
We model intrinsic motivation as maximization of empowerment \citep{klyubin2005empowerment, salge2014empowerment}, defined as the mutual information between an agent's sequence of actions and the subsequent future states of the world. Informally, a highly empowered agent has the potential to have a large effect on the world. An egoistic agent then chooses actions that maximize its own empowerment. An altruistic agent instead chooses actions that maximize the empowerment of another agent. This makes empowerment a natural framework for studying how the direction of intrinsic motivation shapes group behavior.

%% what we do
Prior work has mostly studied empowerment in single-agent settings, leaving its behavior under multi-agent interaction largely unexplored (though see \citealt{capdepuy2012perception, capdepuy2007maximization, salge2017empowerment, charlesworth_flocking}). Recently, \citet{tiomkin_2024} proposed a method for calculating empowerment for arbitrary durations of action sequences and future world sequences using a weak-control approximation around the system's autonomous dynamics, and mapped the problem onto Shannon's classical Gaussian channel. Here we extend these results to the multi-agent case by modeling coupled dynamics as a multi-user interference channel \citep{carleial_1978}. In our proposed approach (Figure~\ref{fig:interference_channel}), each agent treats the actions of others as structured interference noise and solves for its optimal strategy via iterative water-filling \citep{yu_2004}, reaching a Nash equilibrium. This yields a tractable approximation of empowerment in multi-agent systems without requiring explicit coordination or communication between agents.

%% high level results
We study this framework in two qualitatively different classes of systems. The first is the simplest case: a pair of mechanically coupled agents, where empowerment maximization produces dominance hierarchies when agents have unequal strength and cooperation when their strengths are comparable. At the other extreme, we study a large group of agents, modeled as a Vicsek-style flock \citep{vicsek_1995}. Here, egoistic empowerment prevents the usual convergence to a shared heading and instead drives the population into two opposing directional bands. Together these results suggest that empowerment maximization, without any explicit social objective, can generate non-trivial group-level behaviors whose character depends sensitively on agents' relative capabilities and the nature of their coupling.

Our main contributions are as follows: (1) formulation of multi-agent empowerment as an information-theoretic interference channel, (2) efficient solution of the resulting optimization via iterative water-filling, and (3) demonstration of the diversity of behaviors emerging via multi-agent empowerment maximization in few- and many-agent systems.

\section{Related Work}

\subsection{Single Agent Empowerment}\label{sec:single-agent}

For a single agent, empowerment \citep{klyubin2005empowerment,
salge2014empowerment} measures how much influence an agent has over
its future: formally, it is the channel capacity between an agent's
open-loop action sequence $U_{0:t-1} \in \mathbb{R}^{t \cdot d_u}$
and the resulting future state $X_t \in \mathbb{R}^{d_x}$ (one can
also generalize to state sequences \citep{tiomkin_2024}).
In general, an agent need not care about its entire future state. A
\emph{sensor} $C \in \mathbb{R}^{d_y \times d_x}$ selects the
components the agent seeks to control, producing the sensed output
$Y_t = CX_t \in \mathbb{R}^{d_y}$. Empowerment is then the channel
capacity between the action sequence and the sensed future state,
maximized over the agent's probing policy $\pi(\cdot | x_0)$:
\begin{equation*}
    \mathcal{E}(x_0) \coloneq \max_{\pi(\cdot |x_0)}
    \mathcal{I}[Y_t; U_{0:t-1} \mid X_0 = x_0].
\end{equation*}
A tractable approximation proceeds in four steps: (1) treat the
action sequence as a small perturbation to the nonlinear autonomous
dynamics and retain only the linear response, (2) assume the probing
policy is Gaussian,
\begin{equation*}
    \pi(U_{0:t-1} | X_0 = x_0) \coloneq \mathcal{N}(\mathbf{0},
    \Sigma_{u_{0:t-1}}),
\end{equation*}
(3) add observation noise $z \sim \mathcal{N}(\mathbf{0}, \Sigma_z)$
to model a realistic noisy sensor (ensuring finite channel capacity),
and (4) apply the water-filling algorithm for a power-constrained
linear Gaussian channel \citep{cover1999elements}. This gives:
\begin{equation*}
    \mathcal{E}(x_0) \approx \max_{\Sigma_{u_{0:t-1}}} \frac{1}{2}
    \ln \big| CF_t \Sigma_{u_{0:t-1}} F_t^\top C^\top +
    \Sigma_z \big|, \quad \textrm{s.t.} \quad
    \mathrm{tr}(\Sigma_{u_{0:t-1}}) \leq P,
\end{equation*}
where $F_t$ is the sensitivity of the future state to the control
sequence along the autonomous trajectory,
\begin{equation*}
    F_t(x_0) \coloneq \frac{\partial x_t}{\partial
    u_{0:t-1}}\bigg|_{u_{0:t-1} = \mathbf{0}} \in \mathbb{R}^{d_x
    \times (t \cdot d_u)},
\end{equation*}
so that $CF_t \in \mathbb{R}^{d_y \times (t \cdot d_u)}$ is the
effective channel from actions to sensed output. Water-filling
distributes power across the eigenmodes of this effective channel.
When $C = \mathbf{I}$ the sensor is the full state and the original
formulation is recovered. This single-agent formulation does not
extend directly to multiple interacting agents; we address this in Section \ref{sec:method}.

% \subsection{Single Agent Empowerment}\label{sec:single-agent}

% For a single agent, empowerment \citep{klyubin2005empowerment, salge2014empowerment} is defined as the channel capacity between an agent's open-loop action sequence $U_{0:t-1} \subset \mathbb{R}^{t \cdot d_u}$ and the resulting final state $X_t \subset \mathbb{R}^{d_x}$ upon application of those actions on the dynamics (one can generalize empowerment to resulting environmental state sequences as well \citep{tiomkin_2024}).
% This quantity corresponds to the maximum entropy of future states
% achievable by the agent's probing policy $\pi(\cdot | x_0)$ and it
% also depends on the initial state $x_0$ of the agent:
% \begin{equation*}
%     \mathcal{E}(x_0) \coloneq \max_{\pi(\cdot |x_0)} \mathcal{I}[X_t; U_{0:t-1} | X_0 =x_0]
% \end{equation*}
% In general, an agent need not care about its entire future state: a sensor $C \in \mathbb{R}^{d_y \times d_x}$ selects the components of the future state the agent seeks to control, producing the observation $Y_t = C X_t$. Empowerment is then defined as the channel capacity between the agent's action sequence and its sensed future state, and depends on the initial state $x_0$:
% \begin{equation*}
%     \mathcal{E}(x_0) \coloneq \max_{\pi(\cdot |x_0)} \mathcal{I}[Y_t; U_{0:t-1} | X_0 =x_0]
% \end{equation*}
% A tractable approximation to calculate this is to (1) consider the action sequence as a small perturbation to the nonlinear autonomous dynamics, retaining only the linear response, (2) assume the probing policy is Gaussian distributed:
% \begin{equation*}
%     \pi(U_{0:t-1} | X_0 = x_0) \coloneq \mathcal{N}(\mathbf{0}, \Sigma_{u_{0:t-1}}),
% \end{equation*} 
% (3) add observation noise $z \sim \mathcal{N}(\mathbf{0}, \Sigma_z)$
% to model a realistic noisy sensor for the end observation (thus
% ensuring finite channel capacity), and (4) apply the water-filling
% algorithm for channel capacity calculation in a power-constrained
% linear Gaussian system \citep{cover1999elements}. This gives:
% \begin{equation*}
%     \mathcal{E}(x_0) \approx \max_{\Sigma_{u_{0:t-1}}} \frac{1}{2} \ln\det C\Sigma_{x_t}C^\top, \quad \textrm{s.t.} \quad \mathrm{tr}(\Sigma_{u_{0:t-1}}) \leq P, \quad \Sigma_{x_t} = F_t \Sigma_{u_{0:t-1}} F_t^\top + \Sigma_z,
% \end{equation*}
% This amounts to maximizing the log-determinant of the future sensor covariance subject to a trace constraint on the probing power; equivalently, water-filling distributes power across the eigenmodes of the effective linearized channel $CF_t$, the partial derivative of the future sensor $y_t$ with respect to the control sequence $u_{0:t-1}$ evaluated along the autonomous trajectory:
% \begin{equation*}
%     CF_t(x_0) \coloneq \frac{\partial y_t}{\partial u_{0:t-1}}\bigg|_{u_{0:t-1} = \mathbf{0}} \in \mathbb{R}^{d_x \times (t \cdot d_u)}.
% \end{equation*}
% This single-agent formulation does not extend directly to multiple interacting agents; we address this next.

\subsection{Intrinsically Motivated Collective Motion}

Collective motion in animal groups has long served as a testbed for
studying how simple local interactions give rise to group-level
structure
\citep{couzin2009collective,cavagna2014bird,ouellette2022physics}.
\cite{reynolds_1987} showed that rich flocking dynamics can emerge
from a few local decision rules, and \cite{vicsek_1995} distilled
this idea into a minimal active-matter model that has since become a
standard benchmark. Informational accounts of such rules, deriving
them from local information maximization rather than prescribing them
by hand, have also been proposed
\citep{salge11:_local_infor_maxim_emerg_flock_behav}.

\cite{charlesworth_flocking} studied collective behavior under Future State Maximization (FSM), where each agent acts to maximize the number of distinct visual futures it can reach over a finite time horizon. Agents perceive neighbors through a simple retina and choose among a small set of discrete actions by explicitly enumerating reachable future percepts along a rollout tree. Cohesion, co-alignment, and collision avoidance, which all are qualitative hallmarks of animal flocking, emerge spontaneously, without being built into the objective. Because explicit enumeration scales poorly, the authors trained a neural network to approximate FSM and reach larger flock sizes.

The principle of FSM overlaps substantially with empowerment
maximization: both objectives reward agents for preserving a large
set of reachable futures. However, the two frameworks diverge in
important respects. First, our method is analytically tractable: by
casting multi-agent empowerment as an interference channel and
solving via iterative water-filling, we obtain a principled solution
without learned approximations. Second, our agents operate directly
in continuous action and state spaces with coupled dynamics, rather
than over discrete visual percepts and a fixed search tree. We will
furthermore find that the precise choice of ``sensors'' --- i.e.,
what the agent observes about its reachable futures --- matters. In
the experiments presented here, the collective behaviors emerging
under FSM and under empowerment differ markedly: FSM drives
consensus alignment and cohesion, whereas egoistic empowerment
produces a bimodal configuration in which agents self-organize into
opposing directional bands.
% DP: As discussed earlier with Tristan, we do not compare like to
% like (and do not have time to do a full sweep to confirm what would
% be the right match, so I am intentionally limiting the claim - we
% need to do that experiment in the future),
% that's why I removed the following: This suggests that while FSM and
% empowerment share a common motivation, the precise formulation and
% interaction structure can lead to fundamentally different
% group-level outcomes.

\subsection{Social Influence as Intrinsic Motivation}

\cite{jaques2019socialinfluenceintrinsicmotivation} introduce social
influence as an intrinsic reward: each agent is rewarded for
maximizing its causal effect on other agents' actions, regardless of
whether this benefits itself or its targets. The reward is estimated
via a counterfactual over the other agents' action distributions and
used to train neural-network policies with deep reinforcement
learning. Our framework differs from this one in the target of the
objective. Each of our agents seeks to maximize control over its own
future states, treating others as structured interference rather than
as targets of influence. Egoistic and altruistic behaviors then arise
as natural choices of whose empowerment to maximize --- a degree of
freedom absent in the social-influence framework. Empowerment could
in principle also quantify influence over others, which would bring
the two frameworks into closer alignment; we leave that direction for
future work. As in the comparison with FSM above, our formulation
remains analytically tractable and avoids learned approximations.

\section{Method}\label{sec:method}

In this work we extend empowerment to multi-agent systems with many
interacting agents (Figure \ref{fig:interference_channel}). A key challenge is that channel capacity, the
quantity underlying empowerment, is generally intractable for
nonlinear dynamics. We propose a multi-agent extension of
linear-Gaussian empowerment by decomposing the coupled linearized
dynamics around the autonomous nonlinear solution into an
interference-channel structure: each agent treats its own actions as
a signal to its own future state, and the other agents' actions as
noise. This formulation inherits standard solution techniques and
guarantees from multi-user information theory, extending the
single-agent construction of Section~\ref{sec:single-agent} to the
coupled case. Each agent's empowerment is then the capacity of its
own single-user channel under this structured interference. The
agents' probing distributions are determined jointly as the Nash
equilibrium of a non-cooperative game, which we compute by (1)
linearizing the coupled dynamics over a finite planning horizon, (2)
solving for the per-agent probing covariances with the iterative
water-filling (IWF) algorithm \citep{yu_2004}, and (3) selecting
per-agent control actions by empowerment maximization.

\begin{figure}
    \centering
    \includegraphics[width=\linewidth]{CoLLAs/Figures/Multi_User_Interference.pdf}
    \caption{Multi-agent system as an interference channel with $N=3$ agents with identity sensor $C^{(n)} = \mathbf{I}$. Each agent ($n$) attempts to transmit its message (action sequence $U_{0:t-1}^{(n)}$) through the noisy channel to communicate with its receiver (the agent's future state $X_t^{(n)}$). The interference noise generated by the other channel users affects each agent's ability to communicate with its receiver.}
    \label{fig:interference_channel}
\end{figure}

\subsection{Multi-Agent Linearization}

We consider a multi-agent system composed of $N$ agents indexed by $n \in \{0, \ldots, N-1\}$. Let $x_t^{(n)}$ denote the state of agent $n$ at time $t$, and $u_t^{(n)}$ its action. We stack the individual components of all agents into joint column vectors:
\begin{equation*}
    \mathbf{x}_t =
    \begin{bmatrix}
        x_t^{(0)} \\ x_t^{(1)} \\ \vdots \\ x_t^{(N-1)}
    \end{bmatrix}, \quad
    \mathbf{u}_t =
    \begin{bmatrix}
        u_t^{(0)} \\ u_t^{(1)} \\ \vdots \\ u_t^{(N-1)}
    \end{bmatrix}, \quad
    \mathbf{u}_{0:t-1} =
    \begin{bmatrix}
        u_{0:t-1}^{(0)} \\ u_{0:t-1}^{(1)} \\ \vdots \\ u_{0:t-1}^{(N-1)}
    \end{bmatrix},
\end{equation*}
where $u_{0:t-1}^{(n)}$ denotes the action sequence of agent $n$ over the planning horizon $0$ to $t-1$, and $\mathbf{u}_{0:t-1}$ is the concatenation of all agents' sequences. The system evolves in discrete time under coupled, generally nonlinear dynamics:
\begin{equation*}
    \mathbf{x}_{t+1} = f(\mathbf{x}_t, \mathbf{u}_t).
\end{equation*}
To reduce the multi-agent empowerment problem to a tractable linear-Gaussian form, we linearize these dynamics around the autonomous (zero-control) trajectory. The sensitivity of the joint state at horizon $t$ to the joint action sequence is the block Jacobian
\begin{equation}\label{eq:block_sensitivity}
    \mathbf{F}_t \coloneq
    \begin{bmatrix}
        F_t^{(0, 0)} & F_t^{(0, 1)} & \cdots & F_t^{(0, N-1)} \\
        F_t^{(1, 0)} & F_t^{(1, 1)} & \cdots & F_t^{(1, N-1)} \\
        \vdots & \vdots & \ddots & \vdots \\
        F_t^{(N-1, 0)} & F_t^{(N-1, 1)} & \cdots & F_t^{(N-1, N-1)}
    \end{bmatrix}, \quad F_t^{(n, m)} \coloneq \frac{\partial x_t^{(n)}}{\partial u_{0:t-1}^{(m)}}\bigg|_{\mathbf{u}_{0:t-1}=\mathbf{0}} \in \mathbb{R}^{d_x \times (t \cdot d_u)}.
\end{equation}
The diagonal blocks $F_t^{(n,n)}$ describe each agent's direct influence on its own future state; the off-diagonal blocks $F_t^{(n,m)}$ with $m \neq n$ capture the influence of agent $m$ on agent $n$'s future state through the coupling. This block structure is the basis for the interference-channel formulation developed next.


Each agent $n$ is equipped with a sensor $C^{(n)} \in \mathbb{R}^{d_y \times d_x}$ that picks out the aspects of its future state the agent seeks to influence, yielding the sensed output $Y_t^{(n)} = C^{(n)} X_t^{(n)} \in \mathbb{R}^{d_y}$. The effective channel from agent $n$'s own actions to its sensed output is then $C^{(n)} F_t^{(n,n)}$, while the interference arriving from any other agent $m \neq n$ passes through $C^{(n)} F_t^{(n,m)}$. These are the multi-agent analogs of the effective channel $CF_t$ from Section~\ref{sec:single-agent}, now split by source agent into a direct (signal) block for $m=n$ and off-diagonal interference blocks for $m \neq n$.


\subsection{Iterative Water-Filling}

The linearized multi-agent system induces an interference channel \citep{carleial_1978} (Figure~\ref{fig:interference_channel}): each agent $n$ transmits through its control sequence to its own sensed output $Y_t^{(n)}$, and because the dynamics are coupled, the controls of any other agent $m \neq n$ also reach $Y_t^{(n)}$ as structured interference through the off-diagonal blocks of $\mathbf{F}_t$. This is in contrast to the multiple-access setting studied by \citet{yu_2004}, in which several transmitters share a single receiver. Because each agent seeks to control only its own future, each has its own receiver, so the interference-channel formulation is the natural one.

The resulting optimization is a non-cooperative game whose Nash equilibria characterize the optimal probing distributions; we solve it with iterative water-filling \citep{yu_2004} (Algorithm~\ref{alg:iwf_n_agent}), which under the sufficient conditions of \cite{wang2025convergenceiterativewaterfillingmultiuser} converges to a unique fixed point from any initialization.

For each agent $n$, write $F_t^{(n)} \coloneq F_t^{(n,n)}$ for the direct-channel block, and let
\begin{equation}\label{eq:noise-cov}
    \Sigma_{z,t}^{(n)} \coloneq S_z^{(n)} + \sum_{m \neq n} C^{(n)} F_t^{(n,m)} S_t^{(m)} F_t^{(n,m)\top} C^{(n)\top}
\end{equation}
be the effective noise covariance seen by agent $n$ in its sensed-output space, combining the fixed observation noise $S_z^{(n)} \in \mathbb{R}^{d_y \times d_y}$ with the interference from the other agents' probing covariances $\{S_t^{(m)}\}_{m \neq n}$. Each agent's probing covariance is subject to the power budget $\mathrm{tr}(S_t^{(n)}) \leq P^{(n)}$. At each iteration, agent $n$ updates its action covariance by solving the single-user Gaussian channel problem
\begin{equation}\label{eq:agent-cov-update}
    S_t^{(n)} = \arg\max_{S \succeq 0}
    \frac{1}{2} \ln \left| C^{(n)} F_t^{(n)} S F_t^{(n)\top} C^{(n)\top} + \Sigma_{z,t}^{(n)} \right|,
    \quad \mathrm{s.t.} \quad \mathrm{tr}(S) \leq P^{(n)},
\end{equation}
whose solution is given by the standard water-filling algorithm \citep{cover1999elements}. After convergence, agent $n$'s empowerment is
\begin{equation}\label{eq:agent-empowerment-calculation}
    \mathcal{E}_t^{(n)} = \frac{1}{2} \ln \left| C^{(n)} F_t^{(n)} S_t^{(n)} F_t^{(n)\top} C^{(n)\top} + \Sigma_{z,t}^{(n)} \right| - \frac{1}{2} \ln \left| \Sigma_{z,t}^{(n)} \right|.
\end{equation}
At the Nash equilibrium, each agent's probing distribution is a best response to the others', maximizing its empowerment $\mathcal{E}_t^{(n)}$ given the other agents' distributions.

\begin{algorithm}%[h]
\caption{Iterative Water-Filling for $N$ Agents}
\label{alg:iwf_n_agent}
\begin{algorithmic}[1]
\REQUIRE Sensitivity matrix $\mathbf{F}_t$, power constraints $\{P^{(n)}\}_{n=0}^{N-1}$, observation noise covariances $\{S_z^{(n)}\}_{n=0}^{N-1}$, \\ and sensors $\{C^{(n)}\}_{n=0}^{N-1}$

\STATE Initialize $\{S_t^{(n)}\}_{n=0}^{N-1}$ arbitrarily with $\mathrm{tr}(S_t^{(n)}) \leq P^{(n)}$
\WHILE{$\|S_t^{(n)} - S_t^{(n),\text{prev}}\| > \epsilon$ for any $n$}
    \FOR{$n = 0,\dots,N-1$}
        \STATE Compute $\Sigma_{z,t}^{(n)}$ (\eqref{eq:noise-cov})
        \STATE Compute $S_t^{(n)}$ (\eqref{eq:agent-cov-update})
    \ENDFOR
\ENDWHILE
\STATE Evaluate $\mathcal{E}_t^{(n)}$ (\eqref{eq:agent-empowerment-calculation})
\RETURN $\{\mathcal{E}_t^{(n)}\}_{n=0}^{N-1}$
\end{algorithmic}
\end{algorithm}

\subsection{Policy for Control}

Once the empowerment of each agent has been computed via iterative waterfilling, actions can be derived for each agent. In our work we explore behavior based on alternative control strategies:
\begin{enumerate}
    \item \textbf{Egoistic}: an agent maximizes its own empowerment
    \item \textbf{Altruistic}: an agent maximizes the empowerment of another agent.
\end{enumerate}

In the egoistic control policy each agent selects an action which increases its own empowerment. The egoistic policy for the $n$'th agent is given in \eqref{eq:egoistic-policy}. In this maximization, the $n$th agent selects an action which maximizes its own empowerment under the assumption that all other agent's actions are held fixed with zero control.

\begin{align}
    \pi_\textrm{egoistic}^{(n)}(\mathbf{x}_t) &= \arg\max_{u_t^{(n)}} \mathcal{E}_t^{(n)}\big(f(\mathbf{x}_t, \mathbf{u}_t)\big) \label{eq:egoistic-policy}\\
    \pi_\textrm{altruistic}^{(n)}(\mathbf{x}_t) &= \arg\max_{u_t^{(n)}} \mathcal{E}_t^{(m)}\big(f(\mathbf{x}_t, \mathbf{u}_t)\big) \label{eq:altruistic-policy}
\end{align}

% \begin{equation}\label{eq:egoistic-policy}
%     \pi_\textrm{egoistic}^{(n)}(\mathbf{x}_t) = \arg\max_{u_t^{n}} \mathcal{E}_t^{(n)}\big(f(\mathbf{x_t}, \mathbf{u}_t)\big)
% \end{equation}

In the altruistic policy, each agent maximizes the empowerment of another explicitly specified agent $m$. Notably, this policy does not maximize agent $n$'s ability to influence agent $m$'s state, but rather, the empowerment of agent $m$ directly. The altruistic policy is shown in \eqref{eq:altruistic-policy}. This policy directly computes the optimal action such that agent $n$ increases the empowerment of agent $m$.
% \begin{equation}\label{eq:altruistic-policy}
%     \pi_\textrm{altruistic}^{(n)}(\mathbf{x}_t) = \arg\max_{u_t^{(n)}} \mathcal{E}_t^{(m)}\big(f(\mathbf{x_t}, \mathbf{u_t})\big)
% \end{equation}
Similar to \eqref{eq:egoistic-policy} the actions of the other agents are assumed to be zero under this maximization.

\begin{algorithm}[h]
\caption{Online Multi-Agent Iterative Water-Filling for Control}
\label{alg:interaction}
\begin{algorithmic}[1]
\REQUIRE Initial joint state $\mathbf{x}_0$, planning horizon $t$, sensor noise covariance $\{S_z^{(n)}\}_{n=0}^{N-1}$, and sensors $\{C^{(n)}\}_{n=0}^{N-1}$

\FOR{$k = 0$ to $T-1$}
    \STATE Linearize the $t$-step dynamics at $\mathbf{x}_k$ and compute $\mathbf{F}_t$
    \STATE Calculate each agent's empowerment $\{\mathcal{E}_t^{(n)}\}_{n=0}^{N-1}$ by algorithm \ref{alg:iwf_n_agent}
    \STATE Select action for each agent $\mathbf{u}_k$ using either \eqref{eq:egoistic-policy} or \eqref{eq:altruistic-policy}
    \STATE $\mathbf{x}_{k+1} \leftarrow f(\mathbf{x}_k, \mathbf{u}_k)$
\ENDFOR
\end{algorithmic}
\end{algorithm}

\section{Results}

We evaluate whether multi-agent empowerment induces meaningful collective behavior in coupled systems. We study two qualitatively distinct environments: a small-scale ($N=2$), mechanically coupled pendulum system and a large-scale ($N=125$) agent Vicsek-style flock with distributed spatially dependent interactions. In the pendulum system, agents are directly and mechanically coupled, so the influence of one agent on another is immediate and explicit. The flocking setting is more subtle: agents interact only with local neighbors, yet we show that non-trivial global structure emerges across agents that are never in direct contact, as information propagates through the network of transient local connections. Across both domains, we show that our proposed method generates interpretable individual and group-level structure under different control policies.

\subsection{Linked Pendulums}

We first study a two-agent linked-pendulum system in which each agent applies torque to only its own hinge, while the pendulum tips are coupled by an elastic tendon (shown by the red bar in Figure \ref{fig:four_side_by_side}). Each agent's sensor $C^{(n)}$ selects its angle, so empowerment reflects each agent's capacity to influence its own angular position. This setting provides a minimal testbed for studying how empowerment depends on relative control power (or strength) of the agents.

\begin{figure}
    \centering

    \begin{subfigure}[t]{0.24\columnwidth}
        \centering
        \includegraphics[width=\linewidth]{CoLLAs/Figures/left_up.png}
        \caption{{\color{blue} Left Up}}
        \label{fig:left-up}
    \end{subfigure}
    \hfill
    \begin{subfigure}[t]{0.24\columnwidth}
        \centering
        \includegraphics[width=\linewidth]{CoLLAs/Figures/right_up.png}
        \caption{{\color{orange} Right Up}}
        \label{fig:right-up}
    \end{subfigure}
    \hfill
    \begin{subfigure}[t]{0.24\columnwidth}
        \centering
        \includegraphics[width=\linewidth]{CoLLAs/Figures/both_up.png}
        \caption{{\color{OliveGreen} Both Up}}
        \label{fig:both-up}
    \end{subfigure}
    \hfill
    \begin{subfigure}[t]{0.24\columnwidth}
        \centering
        \includegraphics[width=\linewidth]{CoLLAs/Figures/both_down.png}
        \caption{{\color{gray} Neither Up}}
        \label{fig:neither-up}
    \end{subfigure}

    \caption{The four potential outcomes in the linked pendulum environment. In these figures agent $n=0$ is the blue agent on the left and agent $n=1$ is the orange agent on the right. One strong agent dominates the other agent in Figure \ref{fig:left-up} and Figure \ref{fig:right-up}. In Figure \ref{fig:both-up} both agents are similarly powerful and able to achieve the upright position together. In Figure \ref{fig:neither-up} neither agent is able to get upright.}
    \label{fig:four_side_by_side}
\end{figure}

Figure \ref{fig:empowerment_heatmaps} shows that the system exhibits distinct behavioral regimes as a function of the agents' torque relative to one another. Under the egoistic policy (Figure \ref{fig:heatmap_egoistic}), the dominant outcome is typically determined by asymmetry in actuation strength: when one agent is substantially stronger than the other, it tends to reach and hold the upright position while mechanically constraining the weaker agent through the tendon. This produces the asymmetric outcomes shown in Figures \ref{fig:left-up} and \ref{fig:right-up}.

When the agents have more comparable torque limits, a cooperative regime can emerge in which both pendulums reach upright simultaneously (Figure \ref{fig:both-up}). This regime is associated with high empowerment for both agents, although the empowerment achieved by each agent in this configuration is generally lower than in strongly asymmetric states where a single agent dominates (see the empowerment plots in Figure \ref{fig:heatmap_egoistic}). These results suggest that in multi-agent systems, empowerment depends on the relative capabilities and interactions between agents: stronger agents can dominate, while balanced or cooperative interactions allow multiple agents to simultaneously achieve high empowerment.

\begin{figure}[t]
    \centering

    \begin{subfigure}[t]{\columnwidth}
        \centering
        \includegraphics[width=\linewidth]{CoLLAs/Figures/h=130-egoistic.png}
        \caption{\textbf{Egoistic}: each agent maximizes its own empowerment. The policy is given by \eqref{eq:egoistic-policy}.}
        \label{fig:heatmap_egoistic}
    \end{subfigure}

    \vspace{0.5em}

    \begin{subfigure}[t]{\columnwidth}
        \centering
        \includegraphics[width=\linewidth]{CoLLAs/Figures/h=130-ave.png}
        \caption{\textbf{Altruistic}: the right agent maximizes the empowerment of the left agent by following the policy \eqref{eq:altruistic-policy}. The left agent still follows the egoistic policy \eqref{eq:egoistic-policy}.}
        \label{fig:heatmap_altruistic}
    \end{subfigure}

    \caption{Heatmaps for the linked pendulum environment under two empowerment objectives. Top: egoistic policy, where each agent maximizes its own empowerment. Bottom: altruistic policy, where the right agent acts to increase the empowerment of the left agent. In both objectives the agents have a planning horizon of $1.3$ seconds with $dt=0.01$. In the heatmaps indicating the outcome of the interaction a pendulum is considered in the ``up'' position if its angle is within $\pm1$ radian of vertical.}
    \label{fig:empowerment_heatmaps}
\end{figure}

We next consider a mixed-objective setting in which the left agent follows the egoistic policy while the right agent follows the altruistic policy and acts to increase the left agent's empowerment. The resulting heatmaps are shown in Figure \ref{fig:heatmap_altruistic}. In this case, the left agent reaches the upright state across a substantially larger region of the torque space, including regimes where it would not be able to self-right under egoistic control alone. In particular, increasing the right agent's available torque expands the set of low-power conditions under which the left agent can still attain high empowerment. Empirically, when both agents follow egoistic policies we did not observe that either agent could achieve an upright position with less than $\sim0.5$ torque.

These experiments demonstrate that multi-agent empowerment is sensitive to both the relative control capabilities of agents and their policies. Modifying the optimization target from egoistic to altruistic empowerment maximization can qualitatively reshape the reachable outcome landscape, creating opportunities for weaker agents to achieve high empowerment states that are unavailable when both agents follow egoistic control.

\subsection{Vicsek Flock}

We next evaluate our method in a distributed multi-agent setting using a Vicsek-style flock \citep{vicsek_1995, Ferretti_2022} with $N=125$ agents. Each agent moves at constant speed and controls its angular acceleration, while the Vicsek passive dynamics biases its heading toward local neighbor alignment. As in the pendulum setting, each agent's sensor $C^{(n)}$ selects its angle, so empowerment reflects each agent's capacity to control its own heading. Unlike the linked-pendulum system, interactions here are local, time-varying, and collectively mediated through the group.


\begin{figure}
    \centering

    % ---------------- Column 1 ----------------
    \begin{subfigure}[t]{0.24\textwidth}
        \centering
        \fbox{\includegraphics[width=\linewidth]{CoLLAs/Figures/frame_1s.jpg}}

        \fbox{\includegraphics[width=\linewidth]{CoLLAs/Figures/vicsek_angle_distribution_t20.png}}
        \caption{$1$ (s)}
        \label{fig:plot1}
    \end{subfigure}
    \hfill
    % ---------------- Column 2 ----------------
    \begin{subfigure}[t]{0.24\textwidth}
        \centering
        \fbox{\includegraphics[width=\linewidth]{CoLLAs/Figures/frame_50s.jpg}}

        \fbox{\includegraphics[width=\linewidth]{CoLLAs/Figures/vicsek_angle_distribution_t1000.png}}
        \caption{$50$ (s)}
        \label{fig:plot2}
    \end{subfigure}
    \hfill
    % ---------------- Column 3 ----------------
    \begin{subfigure}[t]{0.24\textwidth}
        \centering
        \fbox{\includegraphics[width=\linewidth]{CoLLAs/Figures/frame_100s.jpg}}

        \fbox{\includegraphics[width=\linewidth]{CoLLAs/Figures/vicsek_angle_distribution_t2000.png}}
        \caption{$100$ (s)}
        \label{fig:plot3}
    \end{subfigure}
    \hfill
    % ---------------- Column 4 ----------------
    \begin{subfigure}[t]{0.24\textwidth}
        \centering
        \fbox{\includegraphics[width=\linewidth]{CoLLAs/Figures/frame_200s.jpg}}

        \fbox{\includegraphics[width=\linewidth]{CoLLAs/Figures/vicsek_angle_distribution_t4000.png}}
        \caption{$200$ (s)}
        \label{fig:plot4}
    \end{subfigure}

    \caption{
    Global structure emerges from egoistic empowerment-driven agents. 
    The flock is comprised of $N=125$ agents. 
    In each column, the top panel shows the spatial configuration of agents at a representative time, where dots denote agent positions and trails indicate recent paths. Color corresponds to the angle of each agent. The bottom panel shows the corresponding distribution of heading angles $\Delta \theta$ relative to the average heading of the flock, illustrating the emergence of directional alignment over time.
    }
    \label{fig:flock-global-structure}
\end{figure}

Figure \ref{fig:flock_stats} compares passive flock dynamics with egoistic empowerment-driven control. Under passive Vicsek dynamics, the average empowerment decreases steadily over time (Figure \ref{fig:flock-empowerment}) as the flock converges toward a highly aligned configuration. Agents effectively lose self-control as they become increasingly aligned with the collective. In contrast, when each agent acts to maximize its own empowerment, the average empowerment in the flock remains substantially higher throughout the simulation.

This difference is mirrored in the flock order parameter defined in \cite{vicsek_1995} (Figure \ref{fig:flock-order}). Passive dynamics rapidly drive the system toward global alignment, with the order parameter approaching $1.0$ indicating complete alignment. By contrast, egoistic empowerment suppresses convergence to a single shared heading and keeps the order parameter near zero. In other words, empowerment-seeking agents resist the loss of individual directional freedom that accompanies consensus.

\begin{figure}[ht]
    \centering

    \begin{subfigure}{0.48\textwidth}
        \centering
        \includegraphics[width=\linewidth]{CoLLAs/Figures/empowerment.png}
        \caption{}
        \label{fig:flock-empowerment}
    \end{subfigure}
    \hfill
    \begin{subfigure}{0.48\textwidth}
        \centering
        \includegraphics[width=\linewidth]{CoLLAs/Figures/order.png}
        \caption{}
        \label{fig:flock-order}
    \end{subfigure}
    
    % \caption{Figure \ref{fig:flock-empowerment} shows a comparison of average empowerment of the flock under the egoistic policy versus the passive dynamics of the system. When each agent attempts to increase its empowerment the average empowerment of the flock increases and remains high. The passive behavior of the Vicsek flock according to \citep{Ferretti_2022} results in a steady decrease in the average flock empowerment over time. Figure \ref{fig:flock-order} demonstrates that passive behavior results in a highly ordered flock as time progresses: agents head in the same direction. Under egoistic behavior the order parameter goes almost to zero.}
    \caption{\ref{fig:flock-empowerment} Average flock empowerment over time under egoistic and passive policies. \ref{fig:flock-order} Flock order parameter over time. Egoistic empowerment sustains high individual control and suppresses consensus alignment, while passive dynamics drive the flock toward complete co-alignment.}
    \label{fig:flock_stats}
\end{figure}

Importantly, low order under the egoistic policy does not imply the absence of structure. Figure \ref{fig:flock-global-structure} shows that the flock self-organizes into coherent large-scale spatial patterns despite remaining globally disordered in heading. Starting from uniformly random headings, the population evolves through transient filamentary and tendril-like formations before settling into a single band configuration that wraps around the periodic domain.

The heading distributions in the lower row of Figure \ref{fig:flock-global-structure} clarify this effect. Rather than converging to a unimodal distribution centered on a common direction, the egoistic policy produces an approximately bimodal heading distribution with two opposing preferred directions. This indicates that the agents maintain high empowerment not by dispersing randomly or heading in a common direction, but by organizing into a structured configuration that maximizes the number of agents moving in opposing directions at the population level.

Together, these results show that egoistic empowerment maximization is sufficient to generate non-trivial large-scale structure in a multi-agent system, without any externally prescribed coordination.


\section{Discussion}

%% opening paragraph. Restating
Multi-agent empowerment maximization, without any extrinsic
coordination objective, is sufficient to produce structured and highly
characteristic 
% interpretable % DP: debatable - the Viczek flock is not so
% transparent, so I modified that
behavior. We show that the specific character of the behavior depends
on the nature of the coupling between agents. In the pendulum system,
varying actuation strength leads to different interpretable outcomes
in terms of domination, collaboration or support. In the Vicsek flock,
egoistic empowerment drives the population into opposing directional
bands which arrange globally, spreading their local organization
globally. We expect other agent sensors to produce different, but
still macroscopically organized structural organizations.
% DP: the following was just tautologically stating what we believe to
% be true about empowerment, anyway, so I changed it to be more
% specific: These outcomes reflect a common underlying principle:
% empowerment-seeking agents organize themselves and others in ways
% that preserve maximal individual control, and that organization is
% shaped by the coupling structure of the system.

%% sell the formulation
Crucially, we formulate multi-agent empowerment as an interference
channel, where each agent's channel capacity problem is isolated and
well-defined by treating the effect of other agents as interference
noise, and thus solved by the IWF algorithm. This provides, for the
first time, a tractable and principled basis for the computation of
collective empowerment in the continuum.
%The key structural insight is that each agent's channel capacity problem is well-defined by treating the effect of other agents as interference noise. This framework opens pathways for future work.

%% computational problems
The most immediate expansion of this work aims scaling it to large
collectives. While the full block sensitivity matrix $\mathbf{F}_t$
scales quadratically in the number of agents, which becomes
prohibitive for large $N$ due to memory constraints, in spatially
distributed systems like the Vicsek flock, most off-diagonal coupling
blocks $F_t^{(n,m)}$ are zero or negligible for non-neighboring
agents. This sparsity will allow reducing the IWF computation to only
the local interaction neighborhood of each agent. A complementary
direction will utilize the optimality condition of IWF for
differentiability: because IWF converges to a fixed point defined by a
stationary condition, implicit differentiation can propagate gradients
without unrolling the iterative procedure. This enables the training
of parameterized dynamics models or policy networks end-to-end using
multi-agent empowerment as an objective.

%% goals and multi-agent empowerment
An important direction is to incorporate goals into multi-agent
intrinsic motivations. Once a group has achieved a
high-empowerment configuration (where agents achieve high individual
control) we hypothesize it is ``ready'' to be more effectively
directed toward an external goal. The altruistic pendulum results hint
at this: a weak agent that cannot self-right under egoistic control
alone can do so when a stronger partner acts altruistically. We will
seek to generalize this and to use empowerment maximization as a
structured preconditioning step to set up goal-directed collective
behavior.


\bibliography{collas2025_conference}
\bibliographystyle{collas2025_conference}

% \appendix
% \section{Appendix}


% \begin{figure}
%     \centering
%     \begin{subfigure}{\linewidth}
%         \centering
%         \includegraphics[width=\linewidth]{CoLLAs/Figures/h=75-egoistic.png}
%         \caption{Planning horizon of $0.75$ seconds.}
%         \label{fig:h=75-egoistic}
%     \end{subfigure}
    
%     \vspace{0.5em}
    
%     \begin{subfigure}{\linewidth}
%         \centering
%         \includegraphics[width=\linewidth]{CoLLAs/Figures/h=125-egoistic.png}
%         \caption{Planning horizon of $1.25$ seconds.}
%         \label{fig:h=125-egoistic}
%     \end{subfigure}
    
%     \vspace{0.5em}
    
%     \begin{subfigure}{\linewidth}
%         \centering
%         \includegraphics[width=\linewidth]{CoLLAs/Figures/h=150-egoistic.png}
%         \caption{Planning horizon of $1.50$ seconds.}
%         \label{fig:h=150-egoistic}
%     \end{subfigure}
    
%     \caption{Outcomes for different planning horizon lengths in the linked pendulum setting. Each agent follows an egoistic policy. For a short planning horizon of $0.75$ seconds shown in figure \ref{fig:h=75-egoistic} the collaborative outcome in figure \ref{fig:both-up} is common. If both agents are strong enough they will both end up in the upright position. As horizon increases to $1.25$ seconds the potential for collaboration decreases as shown in figure \ref{fig:h=125-egoistic}, however there is a small zone where if both agents have similar strengths they can end up in the upright position. Finally, once the planning horizon reaches $1.50$ seconds collaboration almost disappears and only the strongest agent will reach the top as shown in figure \ref{fig:h=150-egoistic}.}
%     \label{fig:pendulum-appendix-heatmaps}
% \end{figure}


\end{document}